\documentclass[12pt, a4paper]{article}

\usepackage[top=2.5cm, bottom=2.5cm, left=3cm, right=3cm]{geometry}
\usepackage{mdframed}
\usepackage[utf8]{inputenc}
\usepackage[T1]{fontenc}
\usepackage[spanish]{babel}
\usepackage{amsmath}
\usepackage{float}
\usepackage[margin=10pt, font=small, labelfont=bf]{caption}
\usepackage{booktabs}
\usepackage{graphicx}
\usepackage{xcolor}
\usepackage{hyperref}
\usepackage{longtable}
\usepackage{comment}
\usepackage{setspace}
\usepackage{titlesec}
\usepackage{fancyhdr}

\setlength{\textfloatsep}{18pt plus 2pt minus 4pt}
\setlength{\floatsep}{12pt plus 2pt minus 2pt}
\setlength{\intextsep}{14pt plus 2pt minus 2pt}
\setstretch{1.15}
\setlength{\parindent}{0pt}
\setlength{\parskip}{8pt}

\titleformat{\section}
  {\large\bfseries}{\thesection.}{0.5em}{}
  [\vspace{2pt}\hrule\vspace{4pt}]
\titleformat{\subsection}
  {\normalsize\bfseries}{\thesubsection.}{0.5em}{}
\titlespacing{\section}{0pt}{18pt}{8pt}
\titlespacing{\subsection}{0pt}{14pt}{6pt}

\hypersetup{colorlinks=true, linkcolor=blue!60!black}

\pagestyle{fancy}
\fancyhf{}
\fancyhead[L]{\footnotesize Centro Javeriano de Competitividad}
\fancyhead[R]{\footnotesize \reportnumber}
\cfoot{\thepage}
\setlength{\headheight}{14pt}
\renewcommand{\headrulewidth}{0.4pt}

\geometry{margin=1in}
\graphicspath{{Paper/figures/}{figures/}}
\newcommand{\tablepath}{Paper/tables/}
\newcommand{\figurelang}{es}
\newcommand{\reportseries}{Informes sobre la Productividad en Colombia}
\newcommand{\reportseriesdesc}{Evidencia para entender y fomentar la productividad}
\newcommand{\reportnumber}{Informe 02}
\newcommand{\reporttitle}{Remuneración Laboral por Trabajador y por Hora en Colombia 2010--2025}
\newcommand{\reportsubtitle}{}
\newcommand{\reportauthorone}{Oliver Pardo}
\newcommand{\reportauthoroneaffil}{Profesor asociado del Departamento de Administración, Pontificia Universidad Javeriana y Director del Centro Javeriano de Competitividad}
\newcommand{\reportauthortwo}{Jorge Orozco}
\newcommand{\reportauthortwoaffil}{Consultor de investigación del Centro Javeriano de Competitividad, Pontificia Universidad Javeriana}
\renewcommand{\contentsname}{Tabla de contenido}

\title{\reporttitle\\\reportsubtitle}
\author{\reportauthorone \and \reportauthortwo}
\date{\today}

\begin{document}

\begin{titlepage}
\thispagestyle{empty}
\vspace*{0.6cm}

\begin{center}

\includegraphics[width=0.58\textwidth]{logo_cjc_palatino_horizontal.png}
\end{center}

\vspace{1.5cm}

\begin{center}
{\Large\bfseries \reportseries\par}
\vspace{0.15cm}
{\large \reportseriesdesc\par}

\vspace{1.4cm}

{\large\bfseries \reportnumber\par}
\vspace{0.45cm}
{\Huge\bfseries \reporttitle\par}
\vspace{0.35cm}
{\Large \reportsubtitle\par}
\vspace{0.65cm}
{\large\bfseries \reportauthorone\par}
{\normalsize \reportauthoroneaffil\par}
\vspace{0.25cm}
{\large\bfseries \reportauthortwo\par}
{\normalsize \reportauthortwoaffil\par}
\end{center}

\vfill

\begin{center}
{\large \today\par}
\end{center}

\end{titlepage}

\begin{mdframed}[
  linewidth=1pt,
  linecolor=gray!60,
  backgroundcolor=gray!6,
  innertopmargin=12pt,
  innerbottommargin=12pt,
  innerleftmargin=14pt,
  innerrightmargin=14pt
]
\section*{Resumen ejecutivo}
\addcontentsline{toc}{section}{Resumen ejecutivo}
\textbf{Este informe analiza la evolución de la remuneración laboral en Colombia entre 2010 y 2025.} El ejercicio se concentra en tres medidas: la remuneración laboral total, la remuneración por trabajador y la remuneración por hora trabajada. La primera muestra el tamaño de la masa de ingresos laborales; la segunda muestra cuánto recibe, en promedio, cada persona ocupada; la tercera muestra cuánto se remunera cada hora trabajada.

\textbf{La remuneración laboral total aumentó a una tasa anual de 3,69\% entre 2010 y 2025.} Ese crecimiento combina un aumento de los ocupados de 2,09\% anual y un aumento de la remuneración por trabajador de 1,56\% anual. En otras palabras, una parte importante del aumento de la remuneración laboral total se explica por la expansión del empleo y no solamente por una mayor remuneración promedio de cada trabajador.

\textbf{La remuneración por hora trabajada creció más rápido que la remuneración por trabajador.} Mientras la remuneración mensual por trabajador pasó de 1,5 a 1,9 millones de pesos constantes de 2025, con un crecimiento anual de 1,56\%, la remuneración por hora pasó de 7,4 a 10,1 miles de pesos constantes de 2025, con un crecimiento anual de 2,03\%. La diferencia se explica por la caída de las horas semanales promedio por trabajador, que pasaron de 46,8 a 43,7.

\textbf{La formalidad sigue marcando una diferencia profunda en los niveles de remuneración.} En 2025, la remuneración mensual por trabajador formal fue de 2,9 millones de pesos constantes de 2025, frente a 1,1 millones en el empleo informal. Por hora trabajada, la remuneración formal fue de 14,5 miles de pesos constantes de 2025, frente a 5,9 miles en la informalidad. La brecha no desapareció durante el periodo.
\end{mdframed}

\newpage
\tableofcontents
\newpage

\section{Introducción}

La productividad laboral importa porque, en el largo plazo, es una de las condiciones para que una economía pueda pagar mejores remuneraciones. Pero la relación entre productividad y remuneración no es automática ni mecánica. Una economía puede producir más por trabajador sin que todos los trabajadores reciban esa mejora de la misma manera. También puede aumentar la remuneración por hora y, al mismo tiempo, reducir el número de horas trabajadas por persona ocupada. Por eso, después de estudiar la evolución del producto por trabajador y del producto por hora trabajada, es importante mirar qué ha pasado con la remuneración laboral.

Este informe hace ese ejercicio para Colombia entre 2010 y 2025. La pregunta es sencilla: cuánto ha aumentado la remuneración laboral por trabajador y cuánto ha aumentado la remuneración por hora trabajada. La distinción es importante porque las dos medidas no responden exactamente a lo mismo. La remuneración por trabajador depende de cuánto se paga cada hora, pero también de cuántas horas trabaja cada persona ocupada. La remuneración por hora trabajada, en cambio, permite observar de manera más directa cuánto se remunera el tiempo efectivamente trabajado.

El informe se concentra en el agregado nacional y en la diferencia entre empleo formal e informal. Esta es una primera mirada. Otros cortes del mercado laboral, como sexo, nivel educativo, región, actividad económica, posición ocupacional y tamaño de empresa, deben estudiarse en informes posteriores. La decisión de separar los informes no es cosmética. El mercado laboral colombiano tiene demasiadas dimensiones relevantes para tratarlas todas a la vez sin perder claridad.

\section{Metodología}

\textbf{La medición utiliza los microdatos armonizados de la GEIH para el periodo 2010--2025.} Como en el Informe 01, se excluye 2020 porque el choque de la pandemia y los cambios operativos de la encuesta hacen que ese año sea difícil de comparar con el resto de la serie. Los ocupados se expanden con los factores de expansión de la encuesta y corresponden al promedio anual implícito en la GEIH.

\textbf{La remuneración por trabajador se construye como una remuneración mensual equivalente.} Para cada ocupado se toma su ingreso laboral real por hora, expresado en pesos constantes de 2025, y se multiplica por las horas semanales trabajadas y por 52/12. De esta manera se obtiene una remuneración mensual comparable a través del tiempo. Luego se suma la remuneración mensual de todos los ocupados usando los factores de expansión y se divide por el número de ocupados.

\textbf{La remuneración por hora trabajada se calcula dividiendo la remuneración laboral total entre el total de horas trabajadas.} Esta definición es importante porque mantiene la consistencia entre las variables: la remuneración por trabajador puede leerse como el producto entre la remuneración por hora trabajada y las horas trabajadas por trabajador. En términos de crecimiento, la remuneración por trabajador puede descomponerse entre el crecimiento de la remuneración por hora y el crecimiento de las horas trabajadas por persona ocupada.

\[
\frac{\text{Remuneración}}{\text{Trabajador}}
=
\frac{\text{Remuneración}}{\text{Hora trabajada}}
\times
\frac{\text{Horas trabajadas}}{\text{Trabajador}}
\]

\textbf{La separación entre empleo formal e informal utiliza el indicador de formalidad armonizado en la base GEIH del proyecto.} En términos generales, los ocupados formales son quienes cumplen la condición de afiliación o cotización a la seguridad social usada para construir la variable de formalidad. Los ocupados informales son quienes no cumplen esa condición. El objetivo no es discutir aquí todas las definiciones posibles de informalidad, sino usar una medida consistente a través del tiempo para comparar niveles y crecimientos de remuneración.

\section{Remuneración agregada}

\begin{table}[H]
\centering
\caption{Remuneración laboral, ocupados, horas y remuneración por hora, 2010--2025}
\label{tab:remuneracion_geih_agregado}
\small
\begin{tabular}{@{}p{8.4cm}rrr@{}}
\toprule
Indicador & 2010 & 2025 & Crec. anual \\
\midrule
Remuneración laboral total mensual (Billones de pesos de 2025) & 25,4 & 43,7 & 3,69\% \\
Ocupados (Millones) & 16,8 & 22,9 & 2,09\% \\
Remuneración por trabajador (Millones de pesos de 2025 al mes) & 1,5 & 1,9 & 1,56\% \\
Horas semanales por trabajador & 46,8 & 43,7 & -0,46\% \\
Remuneración por hora trabajada (Miles de pesos de 2025) & 7,4 & 10,1 & 2,03\% \\
\bottomrule
\end{tabular}
\caption*{\footnotesize Nota: La remuneración por trabajador corresponde a una remuneración mensual equivalente en pesos constantes de 2025. La remuneración por hora se calcula como la remuneración laboral total dividida por las horas trabajadas. Fuente: cálculos propios con GEIH del DANE.}
\end{table}


\textbf{La remuneración laboral total creció por la combinación de más ocupados y una mayor remuneración por trabajador.} En pesos constantes de 2025, la remuneración laboral total mensual pasó de 25,4 billones en 2010 a 43,7 billones en 2025, equivalente a un crecimiento anual de 3,69\%. Ese aumento se explica por dos fuerzas: el número de ocupados creció de 16,8 a 22,9 millones, a una tasa anual de 2,09\%, y la remuneración mensual por trabajador aumentó de 1,5 a 1,9 millones, a una tasa anual de 1,56\%. Dicho de otra forma, una parte importante del crecimiento de la masa salarial provino de tener más personas ocupadas.

\textbf{La remuneración por hora trabajada creció más rápido que la remuneración por trabajador.} Entre 2010 y 2025, la remuneración por hora trabajada aumentó de 7,4 a 10,1 miles de pesos constantes de 2025, con un crecimiento anual de 2,03\%. La remuneración por trabajador aumentó menos, 1,56\% anual, porque las horas semanales promedio por trabajador cayeron de 46,8 a 43,7. Esto significa que parte del aumento de la remuneración por hora se tradujo en una mayor remuneración mensual y parte en menos horas trabajadas por persona ocupada.

\begin{figure}[H]
  \centering
  \includegraphics[width=0.95\textwidth]{Paper/figures/fig_remuneracion_geih_descomposicion.png}
  \caption{Descomposición del crecimiento de la remuneración por trabajador, 2010--2025}
  \label{fig:remuneracion_geih_descomposicion}
  \caption*{\footnotesize Nota: las barras muestran crecimientos anualizados. La descomposición debe interpretarse como una aproximación en puntos porcentuales. Fuente: cálculos propios con GEIH del DANE.}
\end{figure}

\textbf{La Figura \ref{fig:remuneracion_geih_descomposicion} resume esta diferencia.} La remuneración por trabajador creció 1,56\% anual. Ese crecimiento combina un aumento de la remuneración por hora de 2,03\% anual y una caída de las horas trabajadas por trabajador de 0,46\% anual. La lectura es parecida a la del Informe 01: mirar solo la remuneración por trabajador deja por fuera una parte de la historia, porque el tiempo trabajado también cambió.

\section{Remuneración formal e informal}

\begin{table}[H]
\centering
\caption{Remuneración laboral por formalidad, 2010--2025}
\label{tab:remuneracion_geih_formalidad}
\footnotesize
\begin{tabular}{@{}p{1.6cm}p{7.0cm}rrr@{}}
\toprule
Grupo & Indicador & 2010 & 2025 & Crec. anual \\
\midrule
Formal & Remuneración laboral total mensual (Billones de pesos de 2025) & 13,7 & 29,3 & 5,20\% \\
Formal & Ocupados (Millones) & 5,4 & 10,1 & 4,33\% \\
Formal & Remuneración por trabajador (Millones de pesos de 2025 al mes) & 2,6 & 2,9 & 0,83\% \\
Formal & Horas semanales por trabajador & 50,0 & 46,1 & -0,54\% \\
Formal & Remuneración por hora trabajada (Miles de pesos de 2025) & 11,8 & 14,5 & 1,38\% \\
Informal & Remuneración laboral total mensual (Billones de pesos de 2025) & 11,2 & 13,4 & 1,22\% \\
Informal & Ocupados (Millones) & 11,2 & 12,5 & 0,71\% \\
Informal & Remuneración por trabajador (Millones de pesos de 2025 al mes) & 1,0 & 1,1 & 0,51\% \\
Informal & Horas semanales por trabajador & 45,5 & 42,0 & -0,53\% \\
Informal & Remuneración por hora trabajada (Miles de pesos de 2025) & 5,1 & 5,9 & 1,05\% \\
\bottomrule
\end{tabular}
\caption*{\footnotesize Nota: La remuneración por trabajador corresponde a una remuneración mensual equivalente en pesos constantes de 2025. La remuneración por hora se calcula como la remuneración laboral total dividida por las horas trabajadas. Fuente: cálculos propios con GEIH del DANE.}
\end{table}


\textbf{El aumento de la remuneración laboral total se concentró en el empleo formal.} La remuneración laboral total mensual de los trabajadores formales creció 5,20\% anual, mientras que la de los trabajadores informales creció 1,22\% anual. Esta diferencia se explica en buena medida por el aumento del número de ocupados formales, que pasó de 5,4 a 10,1 millones. En contraste, el número de ocupados informales aumentó mucho menos, de 11,2 a 12,5 millones.

\textbf{La brecha de niveles entre formales e informales sigue siendo amplia.} En 2025, la remuneración mensual por trabajador formal fue de 2,9 millones de pesos constantes de 2025, frente a 1,1 millones en el empleo informal. La brecha también aparece al mirar la remuneración por hora: 14,5 miles de pesos en el empleo formal frente a 5,9 miles en el empleo informal. Así, una hora trabajada en el empleo formal recibió en 2025 una remuneración de más del doble de la observada en la informalidad.

\textbf{Las remuneraciones crecieron en ambos segmentos, pero el crecimiento fue mayor en el empleo formal.} La remuneración por trabajador formal creció 0,83\% anual, frente a 0,51\% en el empleo informal. Por hora trabajada, la diferencia también favorece al empleo formal: 1,38\% anual frente a 1,05\%. La reducción de las horas semanales fue muy similar en ambos grupos, cercana a 0,5\% anual. Por esa razón, la diferencia entre la remuneración por hora y la remuneración por trabajador aparece tanto en el empleo formal como en el informal.

\begin{figure}[H]
  \centering
  \includegraphics[width=1.0\textwidth]{Paper/figures/fig_remuneracion_geih_indices.png}
  \caption{Remuneración por trabajador y por hora trabajada según formalidad, 2010--2025}
  \label{fig:remuneracion_geih_indices}
  \caption*{\footnotesize Nota: índices con base 2010 = 100. La serie excluye 2020. Fuente: cálculos propios con GEIH del DANE.}
\end{figure}

\textbf{La Figura \ref{fig:remuneracion_geih_indices} muestra que la recuperación posterior a la pandemia fue especialmente fuerte en la remuneración por hora.} En el agregado, la remuneración por hora trabajada alcanzó en 2025 un índice de 135,1, mientras que la remuneración por trabajador llegó a 126,1. En el empleo formal, la remuneración por hora terminó 22,8\% por encima de su nivel de 2010; en el informal, 16,9\%. Sin embargo, en términos de remuneración por trabajador, el avance fue más moderado: 13,2\% en el empleo formal y 7,9\% en el informal.

\section{Conclusiones}

\textbf{La remuneración laboral en Colombia aumentó entre 2010 y 2025, pero el crecimiento no debe leerse solo a partir de la remuneración por trabajador.} La remuneración por hora trabajada creció más rápido porque las horas semanales por trabajador cayeron durante el periodo. Esta diferencia es importante: parte de la mejora en la remuneración por hora no se tradujo en una remuneración mensual proporcionalmente mayor, sino en una menor intensidad horaria.

\textbf{El empleo formal explica una parte central de la mejora agregada.} Los ocupados formales crecieron mucho más rápido que los informales, y sus remuneraciones también aumentaron a mayor ritmo. Esto hizo que la remuneración laboral total formal ganara peso dentro del mercado laboral. Aun así, la informalidad sigue concentrando un número muy alto de trabajadores y mantiene niveles de remuneración muy inferiores.

\textbf{La agenda de productividad y la agenda laboral no deben discutirse por separado.} Aumentar de manera sostenida la remuneración por trabajador exige aumentar la remuneración por hora trabajada, y eso depende de la capacidad de la economía para elevar la productividad del trabajo. Pero también exige entender por qué la informalidad sigue teniendo niveles de remuneración tan bajos y por qué la mejora agregada no se distribuye con la misma fuerza entre todos los trabajadores. Ese diagnóstico debe ser parte de la agenda económica del país.


\end{document}
